\section{Many Mansions: Inequality of Glory Without Envy}

If every soul in heaven beholds the same infinite Essence of God, is every saint equally blessed? Sacred Scripture, the councils, and the unanimous teaching of the Fathers and Doctors answer plainly: in the one house of the Father, glory is unequally distributed according to merit and charity, yet without a trace of envy or lack.

Our Lord announces the distinction on the night before His Passion: ``In my Father's house there are many mansions'' (John 14:2). St.\ Paul lifts our eyes to the night sky: ``One is the glory of the sun, another the glory of the moon, and another the glory of the stars. For star differeth from star in glory. So also is the resurrection of the dead'' (1 Cor 15:41--42). In the parables of the talents and the pounds, the servant who gains five receives five cities and the one who gains ten receives ten (Matt 25:14--29; Luke 19:16--19); ``every man shall receive his own reward, according to his own labour'' (1 Cor 3:8), and ``he who soweth sparingly shall also reap sparingly: and he who soweth in blessings shall also reap blessings'' (2 Cor 9:6).

In the second century, St.\ Irenaeus of Lyons hands down this teaching from ``the presbyters, the disciples of the apostles.'' Linking John 14:2 with the thirtyfold, sixtyfold, and hundredfold yield of the sower (Matt 13:8), he writes: ``and that was on this account the Lord declared, `In My Father's house are many mansions.' For all things belong to God, who supplies all with a suitable dwelling-place; even as His Word says, that a share is allotted to all by the Father, according as each person is or shall be worthy.''\endnote{St.\ Irenaeus of Lyons, \textit{Against Heresies} V.36.2, trans.\ Alexander Roberts and W.\ H.\ Rambaut, ANF vol.\ 1.}

When fourth-century rigorists such as Jovinian leveled all heavenly rewards, the Church resisted at once and later defined the inequality of glory as dogma. In \textit{Laetentur caeli} (6 July 1439), the Council of Florence defined that the blessed clearly behold the Triune God as He is, yet one more perfectly than another according to the diversity of their merits (\textit{pro meritorum diversitate}).\endnote{Council of Florence, \textit{Laetentur caeli} (6 July 1439), DS 1305.} Likewise, in Session VI, Canon 32 (13 January 1547), the Council of Trent defined that the good works of a justified believer, performed in God's grace and Christ's merit, truly merit increase of grace, eternal life, and an increase of heavenly glory.\endnote{Council of Trent, Session VI (13 January 1547), Decree on Justification, Canon 32, DS 1582.} Every hidden act of charity performed in grace enlarges the soul's eternal capacity to know and enjoy God.

St.\ Thomas Aquinas explains how two souls can look upon the same indivisible God and yet one behold Him more deeply than the other (\textit{ST} I, q.\ 12, a.\ 6). The difference lies neither in the Object seen nor in native cleverness, as though a learned philosopher must sit higher than an unlettered saint. It lies solely in the unequal measure of the \textit{lumen gloriae}, and that measure follows charity alone:

\begin{quote}
Hence the intellect which has more of the light of glory will see God the more perfectly; and he will have a fuller participation of the light of glory who has more charity; because where there is the greater charity, there is the more desire; and desire in a certain degree makes the one desiring apt and prepared to receive the object desired. Hence he who possesses the more charity, will see God the more perfectly, and will be the more beatified.\endnote{St.\ Thomas Aquinas, \textit{Summa Theologiae} I, q.\ 12, a.\ 6.}
\end{quote}

The \textit{Supplement} (q.\ 93, aa.\ 2--3) sharpens the metaphor of ``mansions'': ``the unity of a house corresponds to the unity of beatitude, which unity is on the part of the object, and the plurality of mansions corresponds to the differences of beatitude on the part of the blessed.''\endnote{\textit{Supplement}, q.\ 93, aa.\ 2--3 (English Dominicans, 1922).} Mansions are distinguished in heaven (and degrees of punishment in hell), though not in purgatory, because purgatory is a passage rather than a final end. And since ``our actions are meritorious, not by the very substance of the action, but only by the habit of virtue with which they are informed,'' charity---the form of every virtue---fixes the depth of clarity and joy (\textit{Supplement}, q.\ 93, a.\ 3).

Yet a question shaped by fallen earthly hierarchies immediately arises: if some saints shine like the sun and others like smaller stars, will not those in lower mansions feel envy or lack? St.\ Augustine answers in two classic texts. Preaching on John 14:2 in Tractate 67 on John, he reconciles the ``many mansions'' with the single penny given equally to all laborers in Matthew 20:

\begin{quote}
All alike have that penny, which the householder orders to be given to all that have wrought in the vineyard, making no distinction therein between those who have labored less and those who have labored more: by which penny, of course, is signified eternal life, whereto no one any longer lives to a different length than others, since in eternity life has no diversity in its measure. But the many mansions point to the different grades of merit in that one eternal life. \dots\ The saints, like the stars in the sky, obtain in the kingdom different mansions of diverse degrees of brightness; but on account of that one penny no one is cut off from the kingdom; and God will be all in all in such a way, that, as God is love, love will bring it about that what is possessed by each will be common to all. For in this way every one really possesses it, when he loves to see in another what he has not himself. There will not, therefore, be any envying amid this diversity of brightness, since in all of them will be reigning the unity of love.\endnote{St.\ Augustine, \textit{Tractates on the Gospel of John} 67.2, trans.\ John Gibb and James Innes, NPNF1 vol.\ 7.}
\end{quote}

Because charity makes another's glory one's own, the splendor of Our Lady, the Apostles, and the Martyrs multiplies the joy of the least soul in heaven. In the closing chapter of \textit{The City of God}, Augustine adds the analogy of the human body:

\begin{quote}
And in that blessed city there shall be this great blessing, that no inferior shall envy any superior, as now the archangels are not envied by the angels, because no one will wish to be what he has not received, though bound in strictest concord with him who has received; as in the body the finger does not seek to be the eye, though both members are harmoniously included in the complete structure of the body. And thus, along with his gift, greater or less, each shall receive this further gift of contentment to desire no more than he has.\endnote{St.\ Augustine, \textit{The City of God} XXII.30, trans.\ Marcus Dods (Edinburgh: T.\ \& T.\ Clark, 1871).}
\end{quote}

Every vessel in the Father's house---small cup or deep chalice---is filled to the brim. Nor can that concord ever shatter as Lucifer's did, for Augustine adds in the next lines of \textit{The City of God} XXII.30 that heaven does not abolish free will, but perfects it: Adam's first freedom was the power not to sin (\textit{posse non peccare}); the blessed share God's own indefectible holiness, the inability to sin (\textit{non posse peccare}), ``all the more truly free, because set free from delight in sinning to take unfailing delight in not sinning.''

Alongside the essential crown of the Beatific Vision---which Scholastic theology calls the \textit{aurea}, shared by all the blessed in varying degrees---the medieval Doctors recognized an accidental crown of special joy, the \textbf{aureole} (\textit{aureola}), awarded for three signal victories over the flesh, the world, and the devil.\endnote{\textit{Supplement}, q.\ 96, aa.\ 1, 5, 6, 7 (English Dominicans, 1922).} Drawn from St.\ Thomas's \textit{Sentences} commentary, Question 96 of the \textit{Supplement} defines the \textit{aureola} as an added accidental joy ``in the works one has done, in that they have the character of a signal victory'' (\textit{Supplement}, q.\ 96, a.\ 1):

\begin{itemize}
\item \textit{Virgins}, who win a lifelong victory over concupiscence by choosing perpetual incorruption of body and mind for the kingdom of heaven (Matt 19:12), following the Lamb whithersoever He goeth (\textit{Supplement}, q.\ 96, a.\ 5).
\item \textit{Martyrs}, who win the supreme victory over the world and bodily death for Christ---since, as Augustine teaches, ``it is not the pain but the cause that makes the martyr,'' whether confessing the Faith or dying rather than violating an infused virtue (\textit{Supplement}, q.\ 96, a.\ 6, corpus and ad 9).
\item \textit{Doctors and Preachers}, who conquer the devil not only in themselves but by driving him from the hearts of others through lawful preaching and teaching of saving truth, fulfilling Daniel's promise that ``they that instruct many to justice'' shall shine ``as stars for all eternity'' (Dan 12:3; \textit{Supplement}, q.\ 96, a.\ 7).
\end{itemize}

Yet the \textit{Supplement} adds an essential safeguard (\textit{Supplement}, q.\ 96, a.\ 1, ad 2; a.\ 6, ad 3): a married mother or hidden penitent without any aureole may possess deeper charity---and thus a brighter essential vision of God---than one who wears a martyr's palm or a doctor's crown. Charity alone fixes the height of the mansion; the aureoles are jewels upon the crown, rejoicing the whole court of heaven.
